% ============================================================================
%  THREAT MODEL -- drop-in section for the ASIA CCS / ACM CCS paper.
%
%  Assumptions made while writing this (check before using):
%    - Placed after Related Work, before Methodology/Preliminaries -- no
%      paper notation (W, b, G, ...) is used here; that is introduced fresh
%      in Methodology.
%    - Organized around LLG's own three-axis taxonomy (access point / mode /
%      observation), not around any of our own prior drafts: it is the one
%      taxonomy already native to the baselines this paper compares against,
%      and its "observation" axis is exactly the aux-free/aux-dependent split
%      Related Work is built on.
%    - \ourmethod{} is the same macro as in VLI_related_work_en.tex; define
%      it once per document, not in both files.
%    - \ref{sec:related}, \ref{sec:method}, \ref{sec:eval} are forward
%      references to sections this file does not define; bind the labels
%      when you place this text in the full paper.
%    - The auxiliary-set instantiation (a stratified --prop fraction of the
%      labeled test set) is stated as a fact about this paper's experiments,
%      not as part of the adversary's definition.
% ============================================================================

\section{Threat Model}
\label{sec:threat}

We adopt the threat model common to the label-distribution inference
attacks surveyed in Section~\ref{sec:related}, following the taxonomy
Wainakh et al.~\cite{wainakh2022llg} use to define it: an \emph{access
point} at which the adversary obtains information, a \emph{mode} in which
it acts, and an \emph{observation} tier describing how much it is given
beyond the update itself.

\paragraph{Access point.}
The adversary is the central server that aggregates client updates under a
standard federated learning protocol. In each round it broadcasts a global
model to a client, the client runs several steps of local SGD on its own
data, and reports an update of that model back for aggregation; the
adversary's access point is exactly this report. We consider this under
three aggregation rules -- FedAvg, FedProx, and FedDyn -- since the attacks
compared in this paper differ in how they read the client-side
regularization each one adds to that report.

\paragraph{Mode.}
The adversary is \emph{honest-but-curious}: it executes the aggregation
protocol exactly as specified, broadcasting genuine models and aggregating
updates faithfully, and does not poison the model, inject gradients, or
impersonate a client. Its analysis is passive and confined to a single
client's report within a single round; it requires no persistence across
rounds, no access to other clients' updates, and no collusion. This is the
mode assumed by every attack we compare.

\paragraph{Observation.}
Every attack we compare knows the global model's architecture and
parameters as broadcast, the client's reported update for the final
linear (classification) layer -- its weight and bias gradients,
accumulated over the client's local steps, but not updates to earlier
layers -- and the batch size and number of local steps, needed only to
know how many samples the recovered counts must sum to. None of them
know the client's raw inputs or labels. Beyond this common ground, they
differ along two independent axes.

\emph{Side data.} The attacks fall into the same two tiers
Section~\ref{sec:related} organizes the baselines by:
\begin{itemize}
 \item \emph{White-box only.} No data beyond the above. LLG, iLRG, ZLG in
 its original form, and \ourmethod{} operate at this tier.
 \item \emph{White-box with auxiliary data.} In addition, the adversary
 holds a small auxiliary dataset, distributionally similar to, but not
 identical with, the target client's local data. LLG$^+$, ZLG$^+$, and
 RLU require this tier to calibrate a statistic they compare the
 client's update against. In our experiments this set is instantiated
 the same way RLU's own evaluation instantiates it: a fixed,
 class-stratified fraction of the server's labeled test set, held out
 from every client's training data.
\end{itemize}
Because this second tier is strictly stronger, a fair comparison must
report which tier each result assumes; we do so throughout, and
Section~\ref{sec:method} examines what the second tier buys an attack as
the attacked model is trained toward convergence.

\emph{Rate knowledge.} RLU alone additionally needs the learning rate
$\eta$ and the aggregation rule's own regularization coefficient (FedAvg's
$\gamma$ or FedProx's $\mu$): its posterior-search step is derived by
unrolling the SGD recursion over the local window, which these two
quantities parameterize. No other attack we compare, \ourmethod{}
included, uses either quantity -- the closed-form estimators in this
family recover counts from the shape of the update alone.

\paragraph{A note on the observed signal.}
We describe the attacked quantity as the client's \emph{accumulated
gradient} over its local steps rather than the \emph{weight delta} the
server literally receives under FedAvg. The two coincide exactly under the
plain-SGD configuration used throughout this paper (zero momentum, no
weight decay on the classification layer): by the SGD update rule, the
reported delta is $-\eta$ times the sum of the per-step gradients, so
recovering one is equivalent to recovering the other up to the known
scalar $\eta$. This is the same idealization RLU's own estimator relies on
when it reads a quantity $O$ as a ratio of weight- to bias-update, and it
is why FedProx's proximal term can be undone with a closed-form
coefficient ($\bar\rho$ in Section~\ref{sec:method}) rather than treated as
unobserved noise. It is not exact under FedDyn, whose optimizer applies an
additional weight-decay-like correction outside the gradient we read; we
flag this as an open simplification rather than claim the equivalence
there.

\paragraph{Goal and success criterion.}
Given one client's report for a window of $m$ local steps (equivalently
$mB$ training samples, where $B$ is the batch size), the adversary's goal
is to recover the number of samples $N_j$ belonging to each class
$j \in \{1,\dots,C\}$ trained on in that window -- the window's \emph{label
distribution} -- without reconstructing the samples themselves. This is
strictly weaker than full data reconstruction and strictly stronger than
merely detecting which classes are present; Section~\ref{sec:eval}
formalizes both the class-existence and instance-count views of success
(cAcc and irec) used throughout this paper.

\paragraph{Out of scope.}
We do not consider reconstructing the client's inputs, which prior work
treats as a separate, downstream attack that a recovered label
distribution can assist (Section~\ref{sec:related}); membership or
property inference beyond the label distribution of the attacked window;
an actively malicious server that deviates from the stated aggregation
rule; Byzantine or colluding clients; or defenses such as secure
aggregation, gradient compression, or differential privacy, which change
what the adversary in this model is given and which we leave to future
work.
